% ID: MAT-TRI-IDN-002
% Titolo: Semplificazione di angoli associati
% Disciplina: Matematica
% Area: Trigonometria
% Argomento: Identita
% Classe: Quarta secondaria di secondo grado
% Difficolta: 3
% Tipo: dimostrazione
% Tag: matematica, trigonometria, identita
% Fonte: verifica_funzioni_goniometriche_anonimizzata
% Provenienza: materiale fornito dal docente; testo anonimizzato; soluzione redatta per Exergo

\begin{esercizio}
Semplifica l'espressione
    \[
    \sin\left(\alpha-\frac{7}{2}\pi\right)\operatorname{sec}(\pi+\alpha)
    -\tan(2\pi-\alpha)\tan\left(\frac{\pi}{2}-\alpha\right).
    \]
\end{esercizio}

\begin{soluzione}
Il dominio dell'espressione iniziale è $\sin\alpha\ne0$ e $\cos\alpha\ne0$.
Per periodicità $\sin(\alpha-7\pi/2)=\cos\alpha$ e $\sec(\pi+\alpha)=-1/\cos\alpha$:
il primo prodotto vale $-1$. Inoltre $\tan(2\pi-\alpha)=-\tan\alpha$ e
$\tan(\pi/2-\alpha)=\cot\alpha$, quindi il secondo prodotto vale $-1$.
L'espressione vale zero nel suo dominio.
\end{soluzione}
